(1) 令球内镜像电荷
\[
q'=-\frac{a}{b}q,\qquad b'=\frac{a^2}{b}\in(0,a)
\]
位于 $z=b'$。对球面上任意点 $|\boldsymbol r|=a$，到真实电荷与镜像电荷的距离满足 $|\boldsymbol r-b\hat z|=(b/a)|\boldsymbol r-b'\hat z|$。因此
\[
\phi_A(\boldsymbol r)=\frac1{4\pi\epsilon_0}
\left(\frac q{|\boldsymbol r-b\hat z|}+\frac{q'}{|\boldsymbol r-b'\hat z|}\right)
\]
在 $r=a$ 上恒为零，又在球外除真实电荷处满足 Laplace 方程，满足无穷远零势条件，由静电唯一性定理是正确外部电势。

(2) 接地球诱导总电荷等于镜像电荷总量 $Q_A=q'=-qa/b$，可由远场电势的单极项或积分法向电场确认。真实点电荷受镜像吸引力沿 $-z$：
\[
F_{A,z}=\frac{qq'}{4\pi\epsilon_0(b-b')^2}
=-\frac{q^2ab}{4\pi\epsilon_0(b^2-a^2)^2}.
\]
保持接地的电荷—导体相互作用能以 $b\to\infty$ 为零点并考虑感应电荷重排，得到
\[
U_A(b)=\frac12q\phi_{\rm image}(b)
=-\frac{q^2a}{8\pi\epsilon_0(b^2-a^2)}.
\]
直接计算 $-dU_A/db=F_{A,z}$，证实能量与力自洽；其中 $\tfrac12$ 不可遗漏。

(3) 孤立中性球必须净诱导电荷为零，因此在球心再添加镜像电荷 $q_0=-q'=qa/b$。该中心镜像在球面贡献恒定电势，不改变球面等势性，又使镜像总量 $q'+q_0=0$。于是
\[
\phi_B=\frac1{4\pi\epsilon_0}
\left[\frac q{|\boldsymbol r-b\hat z|}+\frac{q'}{|\boldsymbol r-b'\hat z|}+\frac{q_0}{r}\right].
\]
额外中心镜像对真实电荷作用为正的径向斥力 $qq_0/(4\pi\epsilon_0b^2)$，因此
\[
F_{B,z}=-\frac{q^2ab}{4\pi\epsilon_0(b^2-a^2)^2}
+\frac{q^2a}{4\pi\epsilon_0b^3}.
\]
对 $b>a$，两项的比值为 $(b^2-a^2)^2/b^4<1$，所以 $F_{B,z}<0$；中性球仍吸引正点电荷，只是吸引力弱于接地球。

(4) 在 $b\gg a$ 下，接地球 $F_{A,z}\sim-q^2a/(4\pi\epsilon_0b^3)$；中性球展开 $\{(1-a^2/b^2)^{-2}-1\}\sim 2a^2/b^2$，得到
\[
F_{B,z}\sim-\frac{q^2a^3}{2\pi\epsilon_0b^5},
\]
体现中性极化物体更快衰减的偶极感应力。$b\to a^+$ 时两种力均由靠近表面的局部感应电荷主导而发散；有限原子尺度使理想球面模型在极近距离失效。镜像仅是为满足外部静电边界而放在导体内部的数学辅助源，不能计算真实导体内的局部物理场或把它独立加入系统动力学。
